\chapter{Conclusion}
\label{chapter:conclusion}

\section{Summary of contributions}
\label{section:conclusion-summary}

This thesis set out to investigate how Secure Boot can be incorporated into space mission bootloaders without compromising the strict reliability and performance constraints of the domain. The work progressed through five interconnected contributions.

A \textbf{lifecycle-based threat model} was developed, grounded in the ENISA Space Threat Landscape taxonomy~\cite{enisa_2025}. Two attacker scenarios were identified as most relevant to the bootloader: a physical adversary with pre-launch hardware access and a remote adversary exploiting the communications and update channels after deployment. These scenarios motivated every design decision that followed.

A \textbf{formal Secure Boot specification} (SAVOIR-GS-002S extension) was produced, using the language and requirement structure of SAVOIR so that the document can be adopted or reviewed by the relevant standards bodies without requiring a translation step. The specification introduces seven security requirements (SEC.01--SEC.07) covering a hardware Root of Trust, digital signature verification of all ASW images, non-bypassability of the verification step, protection against classical and post-quantum adversaries, rollback protection via a monotonic counter, and structured failure reporting through existing telemetry channels. All seven requirements are additive to SAVOIR-GS-002 and can be adopted without replacing any existing requirements, with the sole exception of a constraint on BEF.22 (Fast Boot Path) that prevents signature verification from being placed in the skippable subset.

A \textbf{hardware trade-off analysis} was conducted to identify the minimal hardware primitives required to realise the specification. Two are strictly necessary: sector-level hardware-enforced write protection for the BSW and public key storage, and protected non-volatile storage for the monotonic rollback counter. Both are present in commodity ARM Cortex-M microcontrollers, such as the STM32F439ZI selected for the proof of concept, which implements them via Flash option bytes. A hardware hash accelerator is beneficial for classical schemes but not required.

A \textbf{proof-of-concept implementation} in C on the STM32F439ZI demonstrated that all seven SEC requirements are achievable on representative COTS hardware. The implementation includes A/B image partitioning, monotonic rollback protection in option-byte-protected non-volatile storage, PUS-compatible boot reporting, and an algorithm-agnostic cryptographic verification interface. An end-to-end automated test suite, communicating with the board over UART and SWD, verified the correct behaviour of all nominal and failure paths.

Finally, \textbf{performance benchmarks} for four signature schemes (ECDSA-P256, RSA-3072/4096, ML-DSA-44/65 and LMS) on the target hardware provided concrete data for mission designers choosing between classical and post-quantum algorithms. The results show that the security overhead introduced by signature verification is modest: even a hybrid ECDSA-P256 plus LMS deployment on a maximum-size 128~kB image completes in approximately 363~ms with hardware SHA-2 acceleration, well within a one-second boot budget.


\section{Answers to research questions}
\label{section:conclusion-RQs}

The central question of this thesis was:

\begin{quote}
How can Secure Boot be incorporated into space mission bootloaders in a way that enhances security without compromising the strict reliability and performance constraints of space systems?
\end{quote}

The answer demonstrated by this work is that Secure Boot can be incorporated by extending the existing SAVOIR-GS-002 framework with a layer of cryptographic verification requirements that are additive to, rather than in conflict with, the existing boot sequence structure. The key insight is that SAVOIR already mandates integrity checks, A/B partitioning and a structured nominal sequence; the security gap is that those checks are CRC-based and provide no authenticity guarantee against a deliberate adversary. Replacing CRC with digital signature verification, anchoring the verification key in hardware-write-protected memory, and adding a monotonic rollback counter in protected non-volatile storage are all that is needed to close this gap. The PoC on the STM32F439ZI confirms that these additions are feasible on widely available COTS hardware and that the resulting boot time overhead is acceptable even in the worst-case hybrid configuration.

\paragraph{How can Secure Boot be adapted using a secure-by-design approach compatible with space mission requirements?}

The specification follows a secure-by-design principle by making security failures non-recoverable within the boot sequence: if signature verification fails, the Boot SW rejects the image and reports the failure; there is no fallback path that permits unsigned code to run. The specification is compatible with space mission requirements because it extends SAVOIR-GS-002 without replacing it~--- existing mission software and ground toolchains continue to operate unchanged, with the only new ground-side obligation being the signing of firmware images before upload. The single exception is a constraint on the Fast Boot Path (BEF.22S) that prevents signature verification from being placed in the skippable subset, which is a minimal and justified restriction.

\paragraph{What is the minimal level of hardware support necessary for implementing Secure Boot in space-grade systems?}

Two hardware primitives are strictly necessary. First, sector-level write protection that is enforced by the Flash controller hardware and cannot be bypassed by software running at any privilege level, used to make the BSW code and the public key immutable at runtime. Second, protected non-volatile storage for the monotonic rollback counter, which must survive resets and must not be writable by the ASW. Both primitives are present in commodity ARM Cortex-M microcontrollers and do not require a dedicated secure enclave, TrustZone, or space-grade radiation-hardened silicon. A hardware hash accelerator improves throughput for classical schemes but is optional; the benchmarks show that all four evaluated schemes remain within a practical boot-time budget in software-only mode.

\paragraph{Which existing terrestrial Secure Boot standards and frameworks can be applied to space environments?}

NIST SP~800-193~\cite{NIST-sp-800-193} provides the most directly applicable normative framework. Its requirements for a hardware Root of Trust, an authenticated update mechanism, non-bypassability, and detection event logging map cleanly onto the SEC requirements defined in this work and are cited as the basis for each. The SUIT firmware update architecture~\cite{rfc9019} informed the image header format, with the header acting as the SUIT manifest and carrying the version identifier, image size, CRC and digital signature. Concepts from Android Verified Boot and wolfBoot~--- A/B partition selection, version-bound rollback protection, and algorithm-agnostic verification APIs~--- transferred well at the design level. Full adoption of UEFI Secure Boot or wolfBoot is not appropriate because their infrastructure complexity and code size are incompatible with the resource and audit requirements of space-grade software.

\paragraph{To what extent should post-quantum cryptographic algorithms be considered and do they necessitate dedicated hardware support?}

Post-quantum algorithms should be considered a mandatory element of the specification for any mission with a planned operational lifetime beyond approximately 2035, the NIST deadline for completing the post-quantum migration. The benchmark results demonstrate that both ML-DSA and LMS are viable on current COTS hardware without dedicated post-quantum hardware acceleration. ML-DSA-44 achieves 22~ms verification in software on a 1~kB image and remains competitive with classical schemes even at 128~kB. LMS requires only 3~kB of working memory and 60~bytes for the public key, making it the most attractive option under extreme memory constraints and the only PQ scheme that BSI and ANSSI currently recommend for standalone use without hybridisation. No dedicated hardware support is needed for either scheme; the existing hardware hash accelerator provides no benefit for ML-DSA (which uses SHAKE-256 internally) and only modest gains for LMS. A hybrid deployment combining a classical scheme with a post-quantum scheme adds latency that is the sum of the two individual verification times; for ECDSA-P256 plus LMS on a 128~kB image with hardware SHA-2, this sum is approximately 363~ms, which is well within a one-second boot budget.


\section{Future work}
\label{section:conclusion-futureWork}

The directions described below were identified during the course of this work but were not pursued, primarily due to time constraints within the scope of a single research project. Some also require hardware that was not available during this work (e.g.\ open-source space-grade LEON processors). They are presented here as concrete next steps for anyone building on this foundation.

\subsection{Modifying other standards to fully support Secure Boot}
\label{subsection:conclusion-futureWork-standards}

The secure boot design presented in this thesis operates at the Boot Software level, but it does not exist in isolation: the ECSS PUS standard \cite{ECSS-E-ST-70-41C} defines a set of ground-commanded memory management operations that can directly interfere with the boot chain if left unauthenticated. Fully supporting secure boot therefore requires that constraints be placed on specific ST[06] telecommands.

\paragraph{Service Type Interface Requirements.}
The following ST[06] telecommands must be subject to authentication and authorisation controls before they can be considered compatible with a secure boot implementation:
\begin{itemize}
    \item \textbf{TC[6,2] / TC[6,11] Load raw memory data area}~--- allows arbitrary data to be written to any absolute address, including the Application Software storage region. An unauthenticated load command can overwrite a valid ASW image with a malicious one.
    \item \textbf{TC[6,1] Load object memory data}~--- equivalent risk for structured-addressed memories; allows overwriting named memory objects without cryptographic verification.
    \item \textbf{TC[6,16] Disable memory write protection}~--- removes hardware write protection from a specified memory region. Issuing this command before TC[6,2] is the standard two-step sequence to overwrite a protected region; without authentication, an attacker can unlock and replace the Boot Software itself.
    \item \textbf{TC[6,19] / TC[6,21] Load by reference}~--- load variants that source data from an on-board reference rather than inline in the telecommand. The referencing mechanism must be validated to prevent redirection to attacker-controlled memory.
\end{itemize}
All of these commands would need to carry an authenticated digest (e.g.~a MAC or digital signature computed over the command parameters and payload) and be protected against replay attacks, for example by including a monotonic sequence number or timestamp verified by the onboard software before execution. Alternatively, they could be limited to specific memory regions such that they cannot interfere with the secure boot process.

\subsection{Trust anchor management and key rotation (RFC 6024, RFC 5280)}
\label{subsection:conclusion-futureWork-publicInfra}

The PoC stores a single public key compiled into the BSW at manufacturing time. This is the minimal viable Root of Trust, but it creates a single point of failure: if the corresponding private key is ever compromised, there is no in-field mechanism to rotate the trust anchor without reprogramming the BSW itself~--- an operation that requires write-protection to be disabled and a new, signed BSW to be uploaded. For a mission spanning a decade or more, this might be an unacceptable operational risk.

A more robust approach, motivated by the Trust Anchor Management requirements in RFC~6024~\cite{rfc6024} and recommended as a key-compromise recovery technique in NIST SP~800-193~\cite{NIST-sp-800-193} (\S3.5.1), introduces a hierarchical key structure. A root key~--- held offline and used rarely~--- authorises a set of subordinate update keys, which are the keys actually used to sign individual firmware images. Update keys can then be rotated after launch by uploading a new key certificate signed by the root, without touching the BSW or the root trust anchor. The revocation mechanism defined in RFC~5280~\cite{rfc5280} (X.509 certificate revocation lists) provides the complementary ability to invalidate a compromised update key before its scheduled expiry. Applying this model to the PoC would require adding a certificate parsing step to the BSW verification sequence and extending the image header to carry the update key certificate alongside the firmware signature, consistent with the layered manifest concept already defined in the SUIT architecture~\cite{rfc9019}.

\subsection{Partial image updates}
\label{subsection:conclusion-futureWork-partialUpdate}

The PoC transmits and verifies a complete signed ASW image for every update, regardless of the size of the change being applied. For missions where the ground-to-spacecraft link bandwidth is limited~--- as is common in deep-space operations~--- this is wasteful: a minor bug fix requires uploading the full 128~kB image rather than only the changed portion.

A segmented update scheme addresses this. The ASW image is divided into self-contained segments, each with its own header carrying a load address and an independent digital signature. An update then consists only of the segments that have changed; the BSW loads each segment into its correct RAM location and verifies its signature independently, leaving unchanged segments untouched in flash. This maps directly onto the component-based image structure defined in the SUIT manifest model~\cite{rfc9019,rfc9124}, which already accommodates partial updates.\n\nThe primary design trade-off is segment size. Since each segment header is dominated by the 4\,kB signature field (sized to accommodate the ML-DSA-65 hybrid case), segments that are too small incur a large per-segment storage overhead; segments that are too large defeat the bandwidth saving. The non-linear relationship between image size and verification latency shown in \autoref{section:eval-benchmarks} adds a further constraint, as the optimal segment size will differ across signature schemes.

\subsection{Secure state management}
\label{subsection:conclusion-futureWork-secureState}

The COMM partition in the PoC is a shared, unprotected flash region used to pass operational status from the ASW back to the BSW across a system reset. Because the COMM partition must be writable by the ASW, the BSW treats it as untrusted and never takes security-critical decisions based solely on the COMM value without first verifying digital signatures. This is a deliberate defence-in-depth decision, but it means that a compromised ASW can write any COMM value it chooses~--- for example, repeatedly signalling a swap request to trigger an iterative rollback, or preventing the BSW from entering standby mode when ground intervention is needed.

Authenticating the COMM partition would close this residual attack path, but no clean on-device solution exists. Generating an authenticated COMM value on-device requires a private key stored on the device, which a compromised ASW would also be able to access, so any scheme that relies on the ASW producing authenticated values is insecure. Schemes based on the ground pre-signing a fixed set of COMM state values at manufacturing time avoid this problem, but introduce a replay vulnerability: unless each signed payload includes a unique freshness token, a compromised ASW can replay any previously valid signed value. Including a sequence number inside the signed payload resolves the replay problem, but requires the ground to know at manufacturing time the exact sequence of COMM state transitions the spacecraft will make during its mission, which is not knowable in advance. Alternatively, freshly ground-signed COMM commands could be sent over the uplink whenever a state change is required, removing the manufacturing-time constraint entirely but making COMM state changes dependent on a live ground contact. The appropriate solution will depend on the mission's link availability and operational tempo; this is left as an open design problem.

\subsection{Hardware security module and TrustZone integration}
\label{subsection:conclusion-futureWork-HWmodule}

RFC 9397 \cite{rfc9397} is an informational document regarding Trusted Execution Environment Provisioning. A Trusted Execution Environment is an environment that enforces that any code within the environment cannot be tampered with, and any data used by such code cannot be read or tampered with by any code outside the environment.

The Root of Trust in the PoC is implemented using Flash option bytes, which provide hardware-enforced write protection for BSW code and the stored public key. Because the BSW and ASW never execute simultaneously~--- the BSW verifies and then transfers control, after which it is dormant~--- this write protection is sufficient to guarantee the integrity of the trust anchor: the ASW cannot overwrite BSW code or the public key. The remaining gap is one of confidentiality rather than integrity: a compromised ASW executing at the same processor privilege level can read BSW memory regions, including the stored public key. Reading the public key does not directly benefit an attacker in the current design (it is a public value by definition), but if private key material were ever moved on-device~--- for example, to support on-device COMM partition authentication or a TEEP-style TA provisioning model~--- such material would be equally readable by a compromised ASW. A TEE addresses this by enforcing a Secure World / Normal World split in hardware: Secure World memory is architecturally inaccessible to Normal World code regardless of privilege level, protecting both the integrity and the confidentiality of sensitive data at runtime, not only at rest.

The PoC was implemented on the STM32F439ZI, which predates the ARM TrustZone-M extension (introduced in Cortex-M23 and Cortex-M33). More recent microcontrollers such as the STM32H5 series include TrustZone-M alongside hardware SHA-3, ECDSA and RSA acceleration, which would allow the BSW and its cryptographic keys to be placed in Secure World memory that is architecturally inaccessible to the ASW. Vorago's next-generation VA5 family (Cortex-M55) also supports TrustZone and targets space-grade applications~\cite{vorago_va4_satellites}.

At the cutting-edge space-grade end, Xiphera has presented a hardware IP core specifically designed to provide post-quantum Secure Boot support for space-grade processors, including the Frontgrade GR765~\cite{xiphera_pq_secure_boot_2024}. This approach offloads cryptographic operations to a dedicated, formally verified IP block, providing stronger isolation and potentially faster verification than a software implementation on the main CPU.

This work deliberately chose the most accessible hardware in the ARM Cortex-M ecosystem to establish that the minimal hardware primitives required for Secure Boot are already present in proven COTS silicon. A natural next step is to validate the design on a TrustZone-enabled microcontroller, assess whether the stronger isolation guarantees justify the additional integration complexity, and repeat the performance benchmarks from \autoref{chapter:eval} on the newer silicon to quantify the benefit of hardware-accelerated post-quantum operations.

\subsection{Validation on LEON/space-grade hardware}
\label{subsection:conclusion-futureWork-LEON}

All PoC benchmarks and functional tests were conducted on the STM32F439ZI. While that platform is representative of the ARM Cortex-M ecosystem and has been used in space flight software, SAVOIR was originally written with LEON/SPARC processors as the primary target~--- specifically the GR712RC, GR740 and the upcoming GR765~--- and the canonical BSW implementation for that architecture is GRBOOT~\cite{GRBOOT}. The framework has since been adopted more broadly: N7 Space has implemented a SAVOIR-compliant BSW on both RISC-V and ARM Cortex-M targets~\cite{N7-Space-BSW}, which is the lineage most closely related to the platform used in this work. Until the SEC requirements are validated on a LEON processor, however, the claim that the specification is compatible with the canonical SAVOIR hardware environment remains unverified.

\subsection{Formal verification of boot chain}
\label{subsection:conclusion-futureWork-formalVerification}

The automated test suite described in \autoref{section:eval-securityAnalysis} provides broad behavioural coverage but cannot prove the absence of bugs: it demonstrates that tested scenarios produce the expected outcome, but subtle state corruption under adversarial timing, unexpected transitions in the option-byte application sequence, or edge cases in the rollback logic that are difficult to express as concrete test inputs may remain undetected.

Formal methods would address this by exhaustively verifying that security-critical properties hold across all reachable states. The BSW state machine is well-suited to model checking: it has a small number of persistent states, clearly defined transitions and a concise set of invariants that must never be violated~--- for example, that unsigned firmware is never transferred to execution, or that the rollback counter is never decremented. A model-checked specification of the state machine could verify these properties formally, providing assurance that no reachable execution path bypasses the security checks. This would also bring the PoC closer to the formal correctness evidence required by ECSS-E-ST-40C for software at higher criticality categories.

\subsection{Runtime detection beyond boot-time verification}
\label{subsection:conclusion-futureWork-runtimeDetection}

The current design provides detection of unauthorised firmware exclusively at boot time: the Boot SW verifies each ASW image before transferring execution. Once the ASW is running, no further integrity checks are performed. NIST SP~800-193 \cite{NIST-sp-800-193} defines a broader \emph{Detection} capability (\S4.3) that includes integrity verification ``at runtime'' in addition to the pre-execution check. For spacecraft, this could mean periodic re-verification of ASW images in flash storage or the use of watchdog-timer-based health checks to detect unexpected corruption due to radiation-induced bit flips or a silent attacker who has modified flash memory without triggering a reboot. Extending the specification to cover a Root of Trust for Detection (RTD) that operates continuously during mission would bring the design closer to full NIST SP~800-193 compliance and increase resilience against in-operation attacks.

\section{Reflection}

\subsection{Research in the broader computer science field}

This work sits at the intersection of three areas that are each undergoing significant change: embedded systems security, post-quantum cryptographic engineering, and the application of terrestrial security practices to space systems. Spacecraft have always communicated with the ground, but modern missions increasingly rely on COTS components, standard networking protocols, and complex software stacks that are far more exposed and documented than the purpose-built systems of early space exploration. This has consistently revealed the same pattern seen in other safety-critical engineering domains: communities that have historically prioritised reliability over security accumulate a security debt that becomes increasingly costly to address as deployments scale and the incentives for attackers grow. The space sector is now beginning to face the same challenge, with the ENISA Space Threat Landscape~\cite{enisa_2025} and ESA's institutional commitment to cybersecurity~\cite{ESA_Tech_Strategy_2023, ESA_Cyber_TEC} marking a turning point analogous to what the industrial control systems community went through after incidents such as Stuxnet and the Ukrainian power grid attacks in the 2010s.

Within the computer science field more broadly, the contribution of this work is to provide concrete, measured evidence that modern cryptographic security mechanisms are deployable on severely resource-constrained embedded platforms. A common assumption in the embedded security literature is that cryptographic verification imposes prohibitive overhead; the benchmark results here show that this assumption does not hold for current hardware and algorithm choices. The implication is that the gap between what is possible and what is deployed in space is primarily an engineering and standardisation problem rather than a fundamental resource constraint.

On the post-quantum side, this work also contributes to an active area of research. The field is working through the implications of NIST's recent post-quantum standardisation~--- including the stateful hash-based schemes LMS and XMSS (SP~800-208) and the PQC competition outputs ML-DSA and SLH-DSA (FIPS~204/205)~--- and a recurring question is whether these schemes are practical on resource-constrained hardware. The results presented here provide evidence for the affirmative: both ML-DSA and LMS are viable without dedicated hardware acceleration, with LMS being particularly attractive for the most constrained deployments. This adds a space-focused measurement to a body of literature that has mostly covered IoT and smart-card platforms.

\subsection{Implications, applicability, and real-world impact}

The most direct stakeholder for this work is ESA and the space engineering community. The specification is intentionally written in the style and language of SAVOIR so that it can be reviewed by the relevant working groups and potentially adopted as a formal extension to SAVOIR-GS-002. If adopted, it would provide mission designers with a clear, normative baseline for Secure Boot without requiring them to navigate the full complexity of terrestrial frameworks such as UEFI or Android Verified Boot, which are incompatible with space constraints. For existing missions based on GRBOOT or N7 Space BSW, the specification also identifies the precise gap~--- the absence of authenticated signature verification~--- that would need to be addressed to achieve the security properties described here.

For satellite operators and ground segment teams, the practical implication is that the attack vectors identified in the ENISA threat landscape~\cite{enisa_2025}~--- firmware corruption and malicious code injection, whether through physical access before launch or through the communications channel after deployment~--- can be closed using COTS hardware without requiring a redesign of the spacecraft's software architecture. The necessary security primitives are already present in mainstream microcontrollers; the main cost is the engineering effort to restructure the boot sequence and establish a key management process on the ground.

There are also limitations that mission designers must be aware of. The security guarantees provided by the design are contingent on the ground-side private key remaining secure. A compromised private key would allow an adversary to sign and install arbitrary firmware, defeating the entire chain of trust. This shifts the security problem from the spacecraft to the ground infrastructure, which is a more familiar and well-studied domain but is not addressed within the scope of this work. Additionally, the design provides boot-time integrity verification only; once the ASW is running, no further cryptographic checks are performed. A sufficiently capable adversary with persistent access to the running system could modify ASW code in RAM without triggering a reboot and therefore without triggering re-verification. Extending the design to cover runtime detection, as noted in \autoref{subsection:conclusion-futureWork-runtimeDetection}, would address this residual risk.

Finally, the choice to evaluate on an ARM Cortex-M4 microcontroller reflects the most widely available and cost-accessible class of space-compatible processors today. The design principles transfer to LEON/SPARC platforms~--- which are the primary targets of SAVOIR~--- but the specific performance numbers, option byte mechanisms, and memory constraints are platform-specific. Validation on a GR processor, as discussed in \autoref{subsection:conclusion-futureWork-LEON}, is therefore an important next step before the specification can be considered verified for its primary intended deployment context.

\section{Use of AI}

Generative AI tools were used during this project in accordance with the TU Delft guidelines on the use of generative AI in MSc end projects. The tools used were Claude (Anthropic) as the primary assistant and ChatGPT (OpenAI) on an occasional basis.

During the research phase, both tools were used to test ideas and explore counterarguments to proposed design decisions. During writing, they were used to correct grammar and spelling, improve sentence clarity, and suggest restructuring of drafts. During implementation, they helped navigate unfamiliar parts of existing codebases (wolfSSL, STM32 HAL), generate boilerplate code, assist with debugging, and produce code documentation.

All AI-generated output was manually reviewed and verified before inclusion. All research ideas, design decisions, arguments, interpretations, and conclusions are the author's own.